\section{结论}\label{sec:conc}
本文提出了面向结构化航拍视觉问答的证据级路由系统，
联合选择传输表示与接收端的回答计算方式。
问题类型规则利用不同问题对证据需求的差异，
逐样本正确性预测与能耗价格则用于调节准确率与能耗之间的权衡。
在 VisDrone 的三种信道模型下，相较于固定图像传输，
类型规则提高了准确率，并将计入的系统能耗降低约 $55\%$。
在 DroneVehicle 上，相较于类型规则，验证集选定的 MLP 工作点
将准确率提高了 $1.67$ 个百分点，同时将计入的系统能耗降低了 $60.1\%$。
学习型路由的收益取决于数据集与接收器，非线性预测并非始终优于线性基线。
这些结果支持根据问题需求同时选择通信内容与接收端计算方式；
主要节能收益来自在检测证据足以回答问题时省去 VLM 推理。
上述收益适用于本文评估的无人机与边缘联合能耗模型，
不意味着无人机端能耗降低，因为机载检测开销可能超过无线传输节省的能耗。
后续工作将在嵌入式平台上评估在线路由的端到端能耗与时延。
